\begin{longtable}{
    %|>{\raggedright\arraybackslash}p{0.27\textwidth}
    |>{\raggedright\arraybackslash}p{0.23\textwidth}
    |p{0.55\textwidth}
    |p{0.32\textwidth}
    |p{0.27\textwidth}|
    } % 5 - 1.35 \ 
\caption{Термины кризиса}
\label{tab:crisis-terms} \\
\LongTableHeader{
    \hline
    \rowcolor{black}
    %\textcolor{white}{\textbf{Русский термин}} &
    \textcolor{white}{\textbf{Термин}} &
    \textcolor{white}{\textbf{Описание}} &
    \textcolor{white}{\textbf{Источник}} &
    \textcolor{white}{\textbf{Цитата из источника}} \\
    \hline
}

\longcaptioneachpage{Источник: составлено автором}

% Кризис
\begin{minipage}[t]{\linewidth}
    \raggedright
    \begin{itemize}
        \item \textbf{кризис (рус.)}
        \item crisis (англ.)
        \item crise (фр.)
        \item handelskrisis / handelskrisen (нем.)
        \item crisi (ит.)
        \item crisis (исп.)
    \end{itemize}
\end{minipage}
&
Слово заимствовано из медицины (греч. \textgreek{κρίνω} – \textquotedbl{}решение», «поворотный момент в болезни\textquotedbl{}); в этом смысле кризис – не фаза, а переломный момент, за которым наступает либо выздоровление, либо гибель. В экономическую речь вошло в XVIII в. и доминировало как ведущий термин с 1825 по 1870-е гг. Описывает острое, внезапное нарушение хозяйственного порядка – нередко ассоциируется с паникой и волной банкротств. После работ Митчелла (1926) начал вытесняться нейтральным «рецессия» в качестве обозначения верхней точки поворота, однако с 1990-х гг. вновь стремительно набирает популярность – в 2000 г. давал около 50\% всех соответствующих заголовков в EconLit.
&
\begin{minipage}[t]{\linewidth}
    \vspace{-7pt}
    \RaggedRight
    \begin{enumerate}
        \item Lemonnier C. (1835). \textit{Crises commerciales. Duckett's Dictionnaire de la conversation et de la lecture}\cite{Lemonnier-1835} (первая специальная статья)
    \end{enumerate}
\end{minipage}
&
\textquotedbl{}A commercial crisis occurs every time that the regularity of the movements of exchange of which trade consists is destroyed, suspended or restrained\textquotedbl{}
\begin{flushright}
    Lemonnier C., 1835\cite{Lemonnier-1835}
\end{flushright}
\\ \hline

% Депрессия
\begin{minipage}[t]{\linewidth}
    \raggedright
    \begin{itemize}
        \item \textbf{депрессия (рус.)}
        \item depression (англ.)
        \item dépression (фр.)
        \item depression (нем., заимств.)
        \item depressione (ит.)
    \end{itemize}
\end{minipage}
&
Медицинская метафора \textquotedbl{}опускания\textquotedbl{}, \textquotedbl{}подавления\textquotedbl{} (ср. барометр настроений); первое отчетливое экономическое противопоставление prosperity / depression – у Смита (1776). В 1870–1890-е гг. стал господствующим термином для обозначения затяжных спадов – в противовес «кризису» как острому событию. В XX в. превратился в фазу цикла, описывающую особенно тяжелую и/или длительную стадию сокращения; первое словарное разграничение crisis / depression – в Bliss (1897). Наиболее известный пример – Великая депрессия 1929–1939 гг.
&
\begin{minipage}[t]{\linewidth}
    \vspace{-7pt}
    \RaggedRight
    \begin{enumerate}
        \item Smith A. (1776). \textit{Wealth of Nations}\cite{Smith-1776} (первый четкий пример)
        \item Bliss W. D. P. (1897). \textit{Encyclopedia of social reform}\cite{Bliss-1897} (первое словарное определение)
    \end{enumerate}
\end{minipage}
&
\textquotedbl{}A time of general difficulty and pressure in commercial and monetary circles, if acute, is called a crisis; if prolonged it is usually called a period of depression\textquotedbl{}
\begin{flushright}
    Bliss W. D. P., 1897\cite{Bliss-1897}
\end{flushright}
\\ \hline

% Паника
\begin{minipage}[t]{\linewidth}
    \raggedright
    \begin{itemize}
        \item \textbf{паника (рус.)}
        \item panic (англ.)
        \item panic (фр., нем., заимств.)
    \end{itemize}
\end{minipage}
&
Восходит к имени Пана – греческого бога, наводившего на людей беспричинный ужас. Первое финансово-экономическое употребление – 1742 г., в контексте последствий системы Джона Ло. Термин описывает внезапный массовый страх, сопровождаемый набегом на банки и распродажей активов: в отличие от \textquotedbl{}кризиса\textquotedbl{}, паника – прежде всего психологическое явление, краткосрочное по природе. Словари последовательно проводят разграничение: паника – финансовая, начинается у спекулянтов; кризис – промышленный, длится дольше. Пережила пик в XIX – нач. XX вв. (паники 1825, 1837, 1847, 1857, 1907), затем вышла из академического оборота, чтобы вернуться вновь после 2008 г.
&
\begin{minipage}[t]{\linewidth}
    \vspace{-7pt}
    \RaggedRight
    \begin{enumerate}
        \item Anonymous Member of Parliament (1742). \textit{Pamphlet on the system of John Law}\cite{Anonymous-1742} (первое финансовое употребление)
        \item Montefiore J. (1803). \textit{Commercial dictionary}\cite{Montefiore-1803} (первое словарное употребление в финансовом значении)
    \end{enumerate}
\end{minipage}
&
\textquotedbl{}The Bank, being unable to stand the Shock of the least Panic, was suddenly broke\textquotedbl{}
\begin{flushright}
    Anonymus, 1742\cite{Anonymous-1742}
\end{flushright}
\\ \hline

% Пузырь
\begin{minipage}[t]{\linewidth}
    \raggedright
    \begin{itemize}
        \item \textbf{пузырь (рус.)}
        \item bubble (англ.)
    \end{itemize}
\end{minipage}
&
Слово восходит к образу мыльного пузыря – быстро растущего и лопающегося. Первое экономическое употребление – сатирическая поэма Свифта о Компании Южных морей (1721, по OED); первая словарная статья – Mortimer (1766). На протяжении XVIII – нач. XIX вв. \textquotedbl{}пузырь\textquotedbl{} означал прежде всего мошенническую схему; в конце XIX в. (Palgrave, 1894) акцент сместился на раздутую стоимость акций. В XX в. термин переосмыслен в рамках теории рациональных ожиданий: \textquotedbl{}пузырь\textquotedbl{} – это рост цен сверх фундаментальной стоимости актива. Описывает специфическое нарушение, предшествующее коллапсу, а не регулярную фазу цикла.
&
\begin{minipage}[t]{\linewidth}
    \vspace{-7pt}
    \RaggedRight
    \begin{enumerate}
        \item Swift J. (1721). \textit{The Bubble}\cite{Swift-1721} (первый пример)
        \item Mortimer T. (1766). \textit{New and complete dictionary of trade and commerce}\cite{Mortimer-1766} (первая специальная словарная статья)
    \end{enumerate}
\end{minipage}
&
\textquotedbl{}Bubbles, by which the public have been tricked and deceived, are of two kinds, viz., 1. Those which we may properly enough term trading-bubbles. And, 2. Stock, or funds-bubbles. The former have been of various kinds, and the latter at different times; as in France and England in 1719 and 1720\textquotedbl{}
\begin{flushright}
    Mortimer T., 1766\cite{Mortimer-1766}
\end{flushright}
\\ \hline

% Рецессия
\begin{minipage}[t]{\linewidth}
    \raggedright
    \begin{itemize}
        \item \textbf{рецессия (рус.)}
        \item recession (англ.)
        \item recession (ит., нем., заимств.)
    \end{itemize}
\end{minipage}
&
Слово заимствовано из медицины (\textquotedbl{}ремиссия болезни\textquotedbl{}). Первое экономическое значение – снижение цен – Bristol Mercury (1826). Как обозначение фазы всей экономики появилось в Iowa Recorder (1905). Митчелловская реформа 1926 г. ввела \textquotedbl{}рецессию\textquotedbl{} как намеренно нейтральную замену \textquotedbl{}кризиса\textquotedbl{} – обозначение верхней точки поворота цикла без импликации \textquotedbl{}органической болезни системы\textquotedbl{}. В современном стандарте (NBER) – период значительного снижения экономической активности, обычно не менее двух кварталов. Термин описывает фазу цикла, а не экстраординарное событие, что и составляет его ключевое идеологическое отличие от \textquotedbl{}кризиса\textquotedbl{}.
&
\begin{minipage}[t]{\linewidth}
    \vspace{-7pt}
    \RaggedRight
    \begin{enumerate}
        \item Anonymous (1826). \textit{Bristol Mercury}\cite{Anonymous-1826} (первое экономическое употребление)
        \item Mitchell W. C. (1927). \textit{Business Cycles: The Problem and Its Setting}\cite{Mitchell-1927} (введение как фазы цикла)
    \end{enumerate}
\end{minipage}
&
\textquotedbl{}Every business cycle includes a phase of recession; this recession may or may not be marked by a crisis\textquotedbl{}
\begin{flushright}
    Mitchell W. C., 1927\cite{Mitchell-1927}
\end{flushright}
\\ \hline

% Стагнация
\begin{minipage}[t]{\linewidth}
    \raggedright
    \begin{itemize}
        \item \textbf{стагнация (рус.)}
        \item stagnation (англ.)
        \item stagnation (фр.)
    \end{itemize}
\end{minipage}
&
Медицинская метафора застоя крови (ср. нем. Blut stagniert – \textquotedbl{}кровь не движется\textquotedbl{}). Первое экономическое употребление – аноним 1724 г. Besomi выделяет пять устойчивых значений: (1) вялость торговли вообще; (2) следствие кризиса во французских словарях; (3) фаза цикла Оверстона (между \textquotedbl{}давлением\textquotedbl{} и \textquotedbl{}бедствием\textquotedbl{}); (4) тезис Хансена о \textquotedbl{}зрелой экономике\textquotedbl{} (secular stagnation, 1938); (5) нижняя фаза длинной волны по Кондратьеву. В XX–XXI вв. термин вновь актуализирован в дискуссиях о \textquotedbl{}вековой стагнации\textquotedbl{} (Summers, 2013).
&
\begin{minipage}[t]{\linewidth}
    \vspace{-7pt}
    \RaggedRight
    \begin{enumerate}
        \item Anonymous (1724). \textit{Pamphlet on trade}\cite{Anonymous-1724} (первое экономическое употребление)
        \item Ganilh C. (1826). \textit{Dictionnaire analytique d'économie politique}\cite{Ganilh-1826} (первая словарная статья)
        \item Loyd S.J. [Lord Overstone] (1837). \textit{Speech on the state of trade}\cite{Loyd-1837} (как фазы цикла)
    \end{enumerate}
\end{minipage}
&
\textquotedbl{}There must either follow an entire Stagnation of Trade, or a proportionable Rise of all foreign Commodities\textquotedbl{}
\begin{flushright}
    Anonymous, 1724\cite{Anonymous-1724}
\end{flushright}
\\ \hline

% Бедствие / Тяготы
\begin{minipage}[t]{\linewidth}
    \raggedright
    \begin{itemize}
        \item \textbf{бедствие / тяготы (рус.)}
        \item distress (англ.)
        \item détresse (фр.)
    \end{itemize}
\end{minipage}
&
Термин описывает прежде всего субъективные последствия кризиса для людей и предприятий, а не само объективное явление – отсюда его определение Besomi как \textquotedbl{}пассивного и субъективного\textquotedbl{}. Применялся к финансовым страданиям с 1720-х гг. (аноним, 1722); приобрел устойчивость с 1793 г. Фигурирует у Оверстона (1837) как последняя фаза десятифазного цикла – после \textquotedbl{}стагнации\textquotedbl{}, перед возвратом к \textquotedbl{}успокоению\textquotedbl{}. К 1840-м гг. вытеснен словами \textquotedbl{}паника\textquotedbl{} и \textquotedbl{}кризис\textquotedbl{}; его место в словарях занял термин \textquotedbl{}депрессия\textquotedbl{}. Французское détresse des finances (Boislandry, 1790) имеет самостоятельную историю.
&
\begin{minipage}[t]{\linewidth}
    \vspace{-7pt}
    \RaggedRight
    \begin{enumerate}
        \item Currie J. (1793). \textit{Pamphlet on commercial distress}\cite{Currie-1793} (первое четкое употребление в значении торгово-финансового бедствия)
        \item Loyd S.J. [Lord Overstone] (1837). \textit{Speech on the state of trade}\cite{Loyd-1837} (как фаза цикла)
    \end{enumerate}
\end{minipage}
&
\textquotedbl{}General distress in the commercial and manufacturing interests\textquotedbl{}
\begin{flushright}
    Currie J., 1793\cite{Currie-1793}
\end{flushright}
\\ \hline

% Затруднение / Стеснение
\begin{minipage}[t]{\linewidth}
    \raggedright
    \begin{itemize}
        \item \textbf{затруднение / стеснение (рус.)}
        \item embarrassment (англ.)
        \item embarras (фр.)
    \end{itemize}
\end{minipage}
&
Слово означает буквально \textquotedbl{}препятствие\textquotedbl{}, \textquotedbl{}помеху»\textquotedbl{}(фр. embarras – запруда, пробка). Первое употребление в смысле «денежного стеснения» – Sinclair (1797). Активно использовалось до 1840-х гг. как обозначение нарушения платежного оборота, цепи коммерческих долгов. Жюгляр (1863) противопоставлял краткосрочные \textquotedbl{}commercial embarrassments\textquotedbl{} (1–2 года) длительным периодам процветания (6–7 лет). Термин описывал состояние затрудненности, а не сам перелом, и потому вел к фаза → дистресс, но не являлся фазой цикла в строгом смысле. Вышел из употребления вместе с \textquotedbl{}distress\textquotedbl{} к середине XIX в.
&
\begin{minipage}[t]{\linewidth}
    \vspace{-7pt}
    \RaggedRight
    \begin{enumerate}
        \item Sinclair J. (1797). \textit{Pamphlet on pecuniary difficulties}\cite{Sinclair-1797} (первое употребление)
        \item Juglar C. (1862). \textit{Des crises commerciales et de leur retour périodique}\cite{Juglar-1862}
    \end{enumerate}
\end{minipage}
&
\textquotedbl{}Whereas commercial embarrassments last a fairly short time, at most one or two years, prosperous times extend over several years, six to seven on average\textquotedbl{}
\begin{flushright}
    Juglar C., 1862\cite{Juglar-1862}
\end{flushright}
\\ \hline

% Затоваривание
\begin{minipage}[t]{\linewidth}
    \raggedright
    \begin{itemize}
        \item \textbf{затоваривание (рус.)}
        \item glut (англ.)
        \item engorgement (фр.)
        \item überfluss (нем.)
    \end{itemize}
\end{minipage}
&
Слово \textquotedbl{}glut\textquotedbl{} (насыщать до отвращения, переполнять) появилось в OED в 1594 г. (Plat); глагол \textquotedbl{}to glut the market\textquotedbl{} – Смит (1624). Центральный термин меркантилистской и классической дискуссии о возможности общего перепроизводства. Мальтус (1827) дал развернутое определение: general glut – такое избыточное предложение разнородных товаров, что все они падают ниже естественной цены без компенсирующего роста цен где-либо еще. Описывает причину или механизм кризиса, а не его фазу в строгом смысле; с победой рикардианской школы термин потерял центральное место в теории, но сохранился как исторический маркер полемики Мальтус–Рикардо.
&
\begin{minipage}[t]{\linewidth}
    \vspace{-7pt}
    \RaggedRight
    \begin{enumerate}
        \item Plat H. (1594). \textit{Early use of \textquotedbl{}glut\textquotedbl{}}\cite{Plat-1594} (первое употребление)
        \item Malthus T. R. (1827). \textit{Definitions in Political Economy}\cite{Malthus-1827} (первое определение)
    \end{enumerate}
\end{minipage}
&
\textquotedbl{}What is meant by a general glut, is such an abundance of a large mass of commodities of different kinds, as to make them all fall below the natural price, or the ordinary costs of production, without any proportionate rise of price in any other equally large mass of commodities\textquotedbl{}
\begin{flushright}
    Malthus T. R., 1827\cite{Malthus-1827}
\end{flushright}
\\ \hline

% Конвульсия / Потрясение
\begin{minipage}[t]{\linewidth}
    \raggedright
    \begin{itemize}
        \item \textbf{конвульсия / потрясение (рус.)}
        \item convulsion (англ.)
    \end{itemize}
\end{minipage}
&
Медицинская метафора судорог, насильственных непроизвольных движений. Фигурирует у Оверстона (1837) как седьмая фаза его десятифазного цикла – между \textquotedbl{}overtrading\textquotedbl{} (чрезмерной торговлей) и \textquotedbl{}pressure\textquotedbl{} (давлением). Besomi помещает термин в категорию «прочих» вместе с revulsion, stagnation, embarrassment, pressure, bubble – то есть он использовался, но не закрепился как аналитический инструмент. Акцентирует насильственный, внезапный характер экономических потрясений; по природе ближе к описанию disturbance (возмущения), нежели регулярной фазы. Употреблялся приблизительно в 1815–1860 гг.
&
\begin{minipage}[t]{\linewidth}
    \vspace{-7pt}
    \RaggedRight
    \begin{enumerate}
        \item Loyd S. J. [Lord Overstone] (1837). \textit{Speech on the state of trade}\cite{Loyd-1837} (как фаза цикла)
        \item Cargill T. (1845). \textit{Pamphlet on cycles of enterprise}\cite{Cargill-1845}
    \end{enumerate}
\end{minipage}
&
\textquotedbl{}First, we find it in a state of quiescence, next improvement, growing confidence, prosperity, excitement, overtrading, convulsion, pressure, stagnation, distress, ending again in quiescence\textquotedbl{}
\begin{flushright}
    Loyd S. J. [Lord Overstone], 1837\cite{Loyd-1837}
\end{flushright}
\\ \hline

% Отлив / Реверсия
\begin{minipage}[t]{\linewidth}
    \raggedright
    \begin{itemize}
        \item \textbf{отлив / реверсия (рус.)}
        \item revulsion (англ.)
    \end{itemize}
\end{minipage}
&
Слово означает буквально «резкий отрыв», «откат» (лат. revellere – тянуть назад). В экономическом контексте описывает внезапный поворот настроений и потоков кредита – из состояния экспансии к свертыванию. Впервые зафиксировано у McCulloch (1819, 1826). Уайт (1882) включил \textquotedbl{}revulsion\textquotedbl{} в последовательность: \textquotedbl{}crisis, another period of commercial depression, revulsion, stagnation». В Morning Chronicle (1855) термин описывал возобновляющийся цикл: \textquotedbl{}after another cycle of prosperity, renewed speculation springs up, a fresh revulsion occurs, a drain of gold follows\textquotedbl{}. Besomi помещает его в группу «прочих» (1815–1860), указывая, что термин не утвердился как аналитическая категория.
&
\begin{minipage}[t]{\linewidth}
    \vspace{-7pt}
    \RaggedRight
    \begin{enumerate}
        \item McCulloch J. R. (1819). \textit{Treatises on commerce}\cite{McCulloch-1819} (первые употребления)
        \item White H. (1882). \textit{Commercial crises. Lalor's Cyclopædia of political science}\cite{White-1882}
    \end{enumerate}
\end{minipage}
&
\textquotedbl{}After another cycle of prosperity, renewed speculation springs up, a fresh revulsion occurs, a drain of gold follows\textquotedbl{}
\begin{flushright}
    Anonymous, 1855\cite{Anonymous-1855}
\end{flushright}
\\ \hline

% Давление / Нажим
\begin{minipage}[t]{\linewidth}
    \raggedright
    \begin{itemize}
        \item \textbf{давление / нажим (рус.)}
        \item pressure (англ.)
        \item pression (фр.)
        \item druck (нем.)
    \end{itemize}
\end{minipage}
&
Термин описывает финансовое напряжение – сжатие кредита, затруднение с рефинансированием долгов. Вошел в устойчивое экономическое употребление в XIX в. как \textquotedbl{}monetary pressure\textquotedbl{} или \textquotedbl{}credit pressure\textquotedbl{}. Фигурирует у Оверстона (1837) как восьмая фаза цикла – непосредственно после \textquotedbl{}конвульсии\textquotedbl{} и перед \textquotedbl{}стагнацией\textquotedbl{}: давление возникает тогда, когда в ответ на начавшийся коллапс участники рынка начинают стремительно сокращать позиции. Besomi включает термин в группу \textquotedbl{}прочих\textquotedbl{} вместе с convulsion, revulsion, stagnation, embarrassment, bubble – он использовался, но не закрепился как самостоятельная аналитическая единица. Современный аналог – \textquotedbl{}credit crunch\textquotedbl{} (кредитное сжатие).
&
\begin{minipage}[t]{\linewidth}
    \vspace{-7pt}
    \RaggedRight
    \begin{enumerate}
        \item Loyd S. J. [Lord Overstone] (1837). \textit{Speech on the state of trade}\cite{Loyd-1837} (как фаза цикла)
    \end{enumerate}
\end{minipage}
&
\textquotedbl{}First, we find it in a state of quiescence, next improvement, growing confidence, prosperity, excitement, overtrading, convulsion, pressure, stagnation, distress, ending again in quiescence\textquotedbl{}
\begin{flushright}
    Loyd S. J. [Lord Overstone], 1837\cite{Loyd-1837}
\end{flushright}
\\ \hline

% Дискредит
\begin{minipage}[t]{\linewidth}
    \raggedright
    \begin{itemize}
        \item \textbf{дискредит (рус.)}
        \item discredit (англ.)
        \item discrédit (фр.)
    \end{itemize}
\end{minipage}
&
Французский термин, обозначающий состояние всеобщей утраты доверия к кредиту и деловым партнерам. Введен Мишелем (Encyclopédie du commerçant, 1839) для характеристики механизма кризиса: \textquotedbl{}those times of general discredit which suspend for some time the regular course of business\textquotedbl{}. Согласно Besomi, термин указывает на \textquotedbl{}suspension of credit and the breaking down of the chain of payments\textquotedbl{} – то есть описывает не фазу цикла, а один из ключевых механизмов его острой фазы. Близкое понятие во французском – état de méfiance générale. Отдельной словарной традиции в немецком и английском не сложилось; в последних эту нишу заняли panic и embarrassment.
&
\begin{minipage}[t]{\linewidth}
    \vspace{-7pt}
    \RaggedRight
    \begin{enumerate}
        \item Michel A. (1839). \textit{Encyclopédie du commerçant}\cite{Michel-1839} (первое развернутое использование)
        \item Courcelle-Seneuil J.-G. (1842). \textit{Dictionnaire politique}\cite{Courcelle-Seneuil-1842}
    \end{enumerate}
\end{minipage}
&
\textquotedbl{}Those times of general discredit which suspend for some time the regular course of business\textquotedbl{}
\begin{flushright}
    Michel A., 1839\cite{Michel-1839}
\end{flushright}
\\ \hline

% Маразм / Вялость
\begin{minipage}[t]{\linewidth}
    \raggedright
    \begin{itemize}
        \item \textbf{вялость (рус.)}
        \item languor / slump (англ.)
        \item marasme (фр.)
    \end{itemize}
\end{minipage}
&
Французский термин, заимствованный из медицины (греч. \textgreek{μαρασμός} – «угасание», «истощение»). Рафалович (1898) противопоставлял острый \textquotedbl{}кризис\textquotedbl{} затяжному \textquotedbl{}вялость\textquotedbl{}: \textquotedbl{}un état de langueur et de marasme prolongé ou de maladie chronique\textquotedbl{}, для которого в английском лучше использовать слово \textquotedbl{}depression\textquotedbl{}. Таким образом, marasme описывает хронически ослабленное состояние экономики – вялость без выздоровления – в отличие от острого кризиса. Во французской экономической литературе применялся параллельно с \textquotedbl{}stagnation\textquotedbl{}; в немецкой и английской традиции отдельного эквивалента не закрепилось. Характеристика \textquotedbl{}хронической болезни\textquotedbl{} сближает его с тезисом о \textquotedbl{}secular stagnation\textquotedbl{} (Hansen, 1938).
&
\begin{minipage}[t]{\linewidth}
    \vspace{-7pt}
    \RaggedRight
    \begin{enumerate}
        \item Raffalovich A. (1898). \textit{Article on marasme}\cite{Raffalovich-1898} (первое зафиксированное экономическое употребление)
        \item Monbrion M. (1838). \textit{Dictionnaire universel du commerce, de la banque et des manufactures}\cite{Monbrion-1838} (смежное употребление)
    \end{enumerate}
\end{minipage}
&
\textquotedbl{}A state of prolonged languor and slump (marasme) or of chronic disease, for which the word 'depression' is better used in English\textquotedbl{}
\begin{flushright}
    Raffalovich A., 1898\cite{Raffalovich-1898}
\end{flushright}
\\ \hline

% Перепроизводство
\begin{minipage}[t]{\linewidth}
    \raggedright
    \begin{itemize}
        \item \textbf{перепроизводство (рус.)}
        \item overproduction (англ.)
        \item surproduction (фр.)
        \item überproduktion (нем.)
    \end{itemize}
\end{minipage}
&
Термин описывает создание товаров в объеме, превышающем платежеспособный спрос, что, по ряду теорий, является структурной причиной кризисов. В немецкой традиции Рошер (1849) разрабатывал понятие produktionskrisen (производственных кризисов) как предпочтительное обозначение периодических потрясений. Маркс и марксистская традиция сделали \textquotedbl{}перепроизводство\textquotedbl{} (überproduktion) центральным объяснительным понятием. В классической политэкономии \textquotedbl{}overproduction\textquotedbl{} было дискуссионным: Мальтус допускал возможность общего перепроизводства (glut), Рикардо и Сэй ее отрицали (закон Сэя). Обозначает причину или механизм кризиса; как самостоятельная фаза цикла не фигурирует, но во многих теориях описывает фазу восходящего бума.
&
\begin{minipage}[t]{\linewidth}
    \vspace{-7pt}
    \RaggedRight
    \begin{enumerate}
        \item Roscher W. (1849). \textit{Die Gegenwart}\cite{Roscher-1849} (первая немецкая словарная статья)
        \item Sismondi J. C. L. (1819). \textit{De la richesse commerciale}\cite{Sismondi-1819} (первое систематическое обсуждение во французской традиции)
        \item Herkner H. (1892). \textit{Krisen. Conrad's Handwörterbuch der Staatswissenschaften}\cite{Herkner-1892}
    \end{enumerate}
\end{minipage}
&
\textquotedbl{}Economic disturbances of various nature – normally disturbances in the equilibrium between production and demand\textquotedbl{}
\begin{flushright}
    Herkner H., 1892\cite{Herkner-1892}
\end{flushright}
\\ \hline

% Крах / Коллапс
\begin{minipage}[t]{\linewidth}
    \raggedright
    \begin{itemize}
        \item \textbf{крах / коллапс (рус.)}
        \item crash / collapse (англ.)
    \end{itemize}
\end{minipage}
&
\textquotedbl{}Crash\textquotedbl{} – звукоподражательное слово, означающее внезапное разрушение; \textquotedbl{}collapse\textquotedbl{} – буквально \textquotedbl{}схлопывание\textquotedbl{}. В экономическом контексте оба термина описывают быстрое и масштабное обрушение цен активов, банковских структур или финансовой системы в целом. Тесно связаны с \textquotedbl{}кризисом\textquotedbl{} и \textquotedbl{}паникой\textquotedbl{}: паника – психологический механизм, коллапс – рыночный результат. В современном употреблении stock market crash (1929, 1987, 2008) – устойчивое сочетание. Как самостоятельная аналитическая фаза цикла не закрепился, оставаясь описательным термином.
&
\begin{minipage}[t]{\linewidth}
    \vspace{-7pt}
    \RaggedRight
    \begin{enumerate}
        \item White H. (1882). \textit{Commercial crises. Lalor's Cyclopædia of political science}\cite{White-1882}
    \end{enumerate}
\end{minipage}
&
\textquotedbl{}Each period of abnormal and exciting prosperity is followed by a violent collapse, whose phases are a money panic, a sudden rise in the rate of interest\textquotedbl{}
\begin{flushright}
    White H., 1882\cite{White-1882}
\end{flushright}
\\ \hline

% Ликвидация
\begin{minipage}[t]{\linewidth}
    \raggedright
    \begin{itemize}
        \item \textbf{ликвидация (рус.)}
        \item liquidation (англ.)
    \end{itemize}
\end{minipage}
&
Термин буквально означает «превращение в жидкость» (лат. liquidus) – то есть конвертацию активов в деньги для погашения долгов. В экономическом контексте \textquotedbl{}liquidation\textquotedbl{} – терминальная фаза острого кризиса или начало выхода из него: принудительные распродажи активов, банкротства, расчет по долгам. Жюгляр использовал это понятие для описания механизма урегулирования кризиса. Шумпетер и австрийская школа рассматривали ликвидацию как необходимую \textquotedbl{}чистку\textquotedbl{} накопленных дисбалансов. В современных теориях финансового ускорителя – ключевой механизм нисходящей спирали.
&
\begin{minipage}[t]{\linewidth}
    \vspace{-7pt}
    \RaggedRight
    \begin{enumerate}
        \item Juglar C. (1862). \textit{Des crises commerciales et de leur retour périodique}\cite{Juglar-1862} (первое теоретически нагруженное употребление)
        \item Mitchell W. C. (1933). \textit{Encyclopaedia of the Social Sciences. Business cycles}\cite{Mitchell-1933}
    \end{enumerate}
\end{minipage}
&
\textquotedbl{}In December, in the midst of the embarrassments of a painful liquidation\textquotedbl{}
\begin{flushright}
    Juglar C., 1862\cite{Juglar-1862}
\end{flushright}
\\ \hline

% Спад / Сжатие
\begin{minipage}[t]{\linewidth}
    \raggedright
    \begin{itemize}
        \item \textbf{спад / сжатие (рус.)}
        \item contraction (англ.)
        \item abschwung (нем.)
    \end{itemize}
\end{minipage}
&
\textquotedbl{}Contraction\textquotedbl{} введен Митчеллом (1926) как нейтральное обозначение нисходящей фазы цикла в рамках четырехфазной схемы: revival → expansion → recession → contraction. \textquotedbl{}Abschwung\textquotedbl{} – немецкий эквивалент, используемый в \textquotedbl{}Konjunktur\textquotedbl{}-литературе 1920–1930-х гг. В отличие от \textquotedbl{}кризиса\textquotedbl{} или \textquotedbl{}депрессии\textquotedbl{}, \textquotedbl{}contraction\textquotedbl{} не несет никаких коннотаций патологии: это технический термин, указывающий лишь на направление движения. Митчелл намеренно выбирал нейтральные слова, чтобы освободить деловой цикл от теоретических предположений.
&
\begin{minipage}[t]{\linewidth}
    \vspace{-7pt}
    \RaggedRight
    \begin{enumerate}
        \item Mitchell W. C. (1933). \textit{Encyclopaedia of the Social Sciences. Business cycles}\cite{Mitchell-1933}
        \item Burns A. F., Mitchell W. C. (1946). \textit{Measuring Business Cycles}\cite{Burns-1946}
    \end{enumerate}
\end{minipage}
&
\textquotedbl{}Business cycles are a type of fluctuation; each cycle includes a phase of revival, expansion, recession and contraction\textquotedbl{}
\begin{flushright}
    Mitchell W. C., 1933\cite{Mitchell-1933}
\end{flushright}
\\ \hline

%\rowcolor{white}
%\longcaption{Источник: составлено автором}
\end{longtable}